\clearpage\begingroup\NoSourceHead\SourceLeaf{329}
\setlength{\fboxsep}{10pt}
\noindent\fbox{\begin{minipage}{\dimexpr\textwidth-2\fboxsep-2\fboxrule\relax}
\fontsize{9}{11}\selectfont
\begin{center}
{\Large CARL FRIEDRICH GAUSS}\\[3pt]
{\LARGE WERKE}\\[7pt]
HERAUSGEGEBEN VON DER GESELLSCHAFT DER WISSENSCHAFTEN\\
ZU GÖTTINGEN\\[4pt]
IN KOMMISSION BEI JULIUS SPRINGER IN BERLIN\\[3pt]
\rule{.22\linewidth}{.3pt}
\end{center}
\noindent Band
\newcommand{\GaussCatalogueItem}[2]{\par\noindent\hangindent=2.4em\hangafter=1\makebox[2.4em][r]{#1\enspace}#2\par}
\GaussCatalogueItem{I.}{DISQUISITIONES ARITHMETICAE. Zweiter Abdruck. 478 Seiten. 1870.\hfill Kart. RM 48.—}
\GaussCatalogueItem{II.}{HÖHERE ARITHMETIK. Zweiter Abdruck. 525 Seiten. 1876.\hfill Kart. RM 53.—}
\GaussCatalogueItem{}{NACHTRAG zum ersten Abdruck des zweiten Bandes. Seite 493—528. 1876.\hfill Kart. RM 3.50}
\GaussCatalogueItem{III.}{ANALYSIS. Zweiter Abdruck. 499 Seiten. 1876.\hfill Kart. RM 50.—}
\GaussCatalogueItem{IV.}{WAHRSCHEINLICHKEITSRECHNUNG UND GEOMETRIE. Zweiter Abdruck. 492 Seiten. 1880.\hfill Kart. RM 50.—}
\GaussCatalogueItem{V}{MATHEMATISCHE PHYSIK. Zweiter Abdruck. 642 Seiten. 1877.\hfill Kart. RM 64.—}
\GaussCatalogueItem{VI.}{ASTRONOMISCHE ABHANDLUNGEN UND AUFSÄTZE. 664 Seiten. 1874.\hfill Kart. RM 69.—}
\GaussCatalogueItem{VII.}{THEORIA MOTUS UND THEORETISCH-ASTRONOMISCHER NACHLASS. (Parabolische Bewegung, Störungen der Ceres und der Pallas, Theorie des Mondes). 650 S. 1906.\hfill Kart. RM 65.—}
\GaussCatalogueItem{VIII.}{ARITHMETIK, ANALYSIS, WAHRSCHEINLICHKEITSRECHNUNG, GEOMETRIE. (Nachträge zu Band I—IV.) 458 Seiten. 1900.\hfill Kart. RM 46.—}
\GaussCatalogueItem{IX.}{GEODÄSIE. (Fortsetzung von Band IV.) 528 Seiten. 1903.\hfill Kart. RM 53.—}
\GaussCatalogueItem{X.}{Abteilung 1. NACHLASS UND BRIEFWECHSEL ZUR REINEN MATHEMATIK. (Nachträge zu Band I—IV und VIII.) TAGEBUCH. 556 Seiten. 1917.\hfill Kart. RM 59.—}
\GaussCatalogueItem{}{NACHBILDUNG DES TAGEBUCHS (Notizenjournals) 1796—1814.\hfill RM 1.20}
\GaussCatalogueItem{XI.}{Abteilung 1. NACHTRÄGE ZUR PHYSIK, CHRONOLOGIE UND ASTRONOMIE. (Nachträge zu Band V—VII.) 518 Seiten. 1927.\hfill Brosch. RM 47. —. Kart. RM 48.—}
\GaussCatalogueItem{XII.}{VARIA. ATLAS DES ERDMAGNETISMUS. Mit Textfiguren, 1 Übersichtstafel, 4 Zahlentafeln und 15 Kartentafeln. V, 415 Seiten. 1929.\hfill Brosch. RM 49.—. Kart. RM 50.—}
\begin{center}\rule{.22\linewidth}{.3pt}\end{center}
\end{minipage}}
\clearpage
\noindent\fbox{\begin{minipage}{\dimexpr\textwidth-2\fboxsep-2\fboxrule\relax}
\fontsize{9}{11}\selectfont
\noindent X\quad Abteilung 2. ABHANDLUNGEN ÜBER GAUSS' WISSENSCHAFTLICHE TÄTIGKEIT AUF DEN GEBIETEN DER REINEN MATHEMATIK.
\par\noindent\hspace*{2em}Abhandlung 1 und 5. ÜBER GAUSS' ZAHLENTHEORETISCHE ARBEITEN. Von \textsc{Paul Bachmann}. — GAUSS UND DIE VARIATIONSRECHNUNG. Von \textsc{Oscar Bolza}. \(75+65\) Seiten. 1922.\hfill Brosch. RM 17 —
\par\noindent\hspace*{2em}Abhandlung 2 und 3: ÜBER GAUSS' FUNKTIONENTHEORETISCHE ARBEITEN. Von \textsc{Ludwig Schlesinger}. — ÜBER DEN ERSTEN UND VIERTEN GAUSSSCHEN BEWEIS DES FUNDAMENTALSATZES DER ALGEBRA. Von \textsc{Alexander Ostrowski}. \(220+18\) Seiten. 1933.\hfill Brosch. RM 27 —
\par\noindent\hspace*{2em}Abhandlung 4: GAUSS ALS GEOMETER. Von \textsc{Paul Staeckel}. 123 Seiten. 1923.\hfill Brosch. RM 12.50
\par\noindent\hspace*{2em}Abhandlung 6: GAUSS ALS ZAHLENRECHNER. Von \textsc{Philipp Maennchen}. 76 Seiten. 1930.\hfill Brosch. RM 8.40
\par\noindent\hspace*{2em}Abhandlung 7 in Vorbereitung.
\par\medskip\noindent XI. Abteilung 2. ABHANDLUNGEN ÜBER GAUSS' WISSENSCHAFTLICHE TÄTIGKEIT AUF DEN GEBIETEN DER ANGEWANDTEN MATHEMATIK (GEODÄSIE, PHYSIK UND ASTRONOMIE). 642 Seiten. 1929.\hfill Brosch. RM 70.—. Kart. RM 71.—
\par\noindent\hspace*{3em}Hieraus einzeln:
\par\noindent\hspace*{2em}Abhandlung 1: ÜBER DIE GEODÄTISCHEN ARBEITEN VON GAUSS. Von \textsc{A. Galle}. 165 Seiten. 1924.\hfill Brosch. RM 17.—
\par\noindent\hspace*{2em}Abhandlung 2: ÜBER GAUSS' PHYSIKALISCHE ARBEITEN (MAGNETISMUS, ELEKTRODYNAMIK, OPTIK). Von \textsc{Clemens Schaefer}. Mit 12 Figuren. 217 Seiten. 1929.\hfill Brosch. RM 24.—
\par\noindent\hspace*{2em}Abhandlung 3: ÜBER DIE ASTRONOMISCHEN ARBEITEN VON GAUSS. Von \textsc{Martin Brendel}. VI, 258 Seiten. 1929.\hfill Brosch RM 29.—
\par\noindent\hspace*{3em}Die Bände X, 2 und XI, 2 bestehen aus einzelnen Abhandlungen (Essays), die besonders paginiert sind und in der Reihenfolge ihrer Fertigstellung auch gesondert ausgegeben werden. Mit der letzten Abhandlung eines jeden Bandes wird Titelblatt und Inhaltsverzeichnis geliefert.
\par\noindent Supplementband. REGISTER. BESCHREIBUNG DES NACHLASSES. BIOGRAPHISCHES.
\par\noindent\hfill In Vorbereitung.
\par\medskip\noindent\hspace*{2em}\textit{Auf die Preise aller vor 1932 erschienenen Bände wird ein Notnachlass von \(10{}^0\!/\!{}_0\) gewährt.}
\par\bigskip\noindent\rule{\linewidth}{.25pt}\par\vspace{-5pt}\noindent\rule{\linewidth}{.25pt}
\begin{center}\fontsize{7}{8.5}\selectfont GÖTTINGEN, GEDRUCKT IN DER DIETERICHSCHEN UNIVERSITÄTS-BUCHDRUCKEREI (W. FR. KAESTNER)\end{center}
\end{minipage}}
\par\smallskip\begin{center}\footnotesize PRINTED IN GERMANY.\end{center}
\clearpage\endgroup
